
\documentclass{article}
\usepackage{colortbl}
\usepackage{makecell}
\usepackage{multirow}
\usepackage{supertabular}

\begin{document}

\newcounter{utterance}

\centering \large Interaction Transcript for game `cladder', experiment `full\_v1.5\_default', episode 7019 with refineonly{-}Qwen3.5{-}4B.
\vspace{24pt}

{ \footnotesize  \setcounter{utterance}{1}
\setlength{\tabcolsep}{0pt}
\begin{supertabular}{c@{$\;$}|p{.15\linewidth}@{}p{.15\linewidth}p{.15\linewidth}p{.15\linewidth}p{.15\linewidth}p{.15\linewidth}}
    \# & \multicolumn{2}{c}{Player} && \multicolumn{2}{c}{Game Master} \\
    \hline

    \theutterance \stepcounter{utterance}  
    & & & \multicolumn{4}{p{0.6\linewidth}}{
        \cellcolor[rgb]{0.9,0.9,0.9}{
            \makecell[{{p{\linewidth}}}]{
                \texttt{\tiny{[P1$\langle$GM]}}
                \texttt{You are an expert in causal inference. The following question is not a typical commonsense query, but rather a meticulously designed question created by a professor specializing in causal inference, intended to assess the students' mastery of the course content.} \\
\texttt{Imagine a self{-}contained, hypothetical world with only the following conditions, and without any unmentioned factors or causal relationships: Hwax has a direct effect on gyzp and jyka. Jyka has a direct effect on lirg. Gyzp has a direct effect on lirg. Method 1: We look at how jyka correlates with lirg case by case according to gyzp. Method 2: We look directly at how jyka correlates with lirg in general. To understand how jyka affects lirg, is it more correct to use the Method 1 than Method 2?} \\
\texttt{Start your answer with 'Yes' or 'No', followed by additional reasoning or evidence to support your explanation.} \\
            }
        }
    }
     \\ \\

    \theutterance \stepcounter{utterance}  
    & & \multicolumn{4}{p{0.6\linewidth}}{
        \cellcolor[rgb]{0.95,0.95,0.95}{
            \makecell[{{p{\linewidth}}}]{
                \texttt{\tiny{[P1$\leftrightarrow$P]}}
                \texttt{{[}DRAFT{]}} \\
\texttt{Yes} \\
\\ 
\texttt{**Reasoning:**} \\
\\ 
\texttt{To determine the correct method for assessing the causal effect of \$jyka\$ on \$lirg\$, we must analyze the causal structure and the presence of confounding variables.} \\
\\ 
\texttt{1.  **Causal Structure Analysis**:} \\
\texttt{    *   \$gyzp \textbackslash{}rightarrow jyka\$} \\
\texttt{    *   \$gyzp \textbackslash{}rightarrow lirg\$} \\
\texttt{    *   \$jyka \textbackslash{}rightarrow lirg\$} \\
\texttt{    *   This structure forms a **collider** at \$lirg\$ if we consider the paths from \$gyzp\$ to \$lirg\$ via both \$jyka\$ and directly. However, more critically, \$gyzp\$ is a **common cause** (confounder) of both \$jyka\$ and \$lirg\$.} \\
\\ 
\texttt{2.  **Evaluation of Method 2 (General Correlation)**:} \\
\texttt{    *   Method 2 looks at the raw correlation between \$jyka\$ and \$lirg\$ without conditioning on \$gyzp\$.} \\
\texttt{    *   Because \$gyzp\$ causes both \$jyka\$ and \$lirg\$, there is a spurious association (confounding bias) between them. Even if \$jyka\$ had no causal effect on \$lirg\$, the correlation would exist solely due to the shared cause \$gyzp\$.} \\
\texttt{    *   Therefore, Method 2 cannot isolate the direct causal effect of \$jyka\$ on \$lirg\$ because it fails to control for the confounder \$gyzp\$.} \\
\\ 
\texttt{3.  **Evaluation of Method 1 (Conditional Correlation)**:} \\
\texttt{    *   Method 1 looks at the correlation between \$jyka\$ and \$lirg\$ *case by case according to* \$gyzp\$. This is equivalent to conditioning on \$gyzp\$ (stratifying by levels of \$gyzp\$).} \\
\texttt{    *   By fixing the value of the confounder \$gyzp\$, we block the backdoor path from \$jyka\$ to \$lirg\$ that runs through \$gyzp\$.} \\
\texttt{    *   Once \$gyzp\$ is controlled for, the remaining correlation between \$jyka\$ and \$lirg\$ reflects only the direct causal effect of \$jyka\$ on \$lirg\$.} \\
\\ 
\texttt{**Conclusion**:} \\
\texttt{Since \$gyzp\$ is a confounder that creates a spurious association between \$jyka\$ and \$lirg\$, the general correlation (Method 2) is biased. Conditioning on the confounder (Method 1) is the standard causal inference procedure (specifically, closing the backdoor path) to obtain an unbiased estimate of the causal effect. Thus, Method 1 is more correct.} \\
            }
        }
    }
    & \\ \\

    \theutterance \stepcounter{utterance}  
    & \multicolumn{4}{p{0.6\linewidth}}{
        \cellcolor[rgb]{0.561,0.737,0.561}{
            \makecell[{{p{\linewidth}}}]{
                \texttt{\tiny{[P1$\rangle$GM]}}
                \texttt{Yes} \\
\\ 
\texttt{**Reasoning:**} \\
\\ 
\texttt{To determine the correct method for assessing the causal effect of \$jyka\$ on \$lirg\$, we must analyze the causal structure and the presence of confounding variables.} \\
\\ 
\texttt{1.  **Causal Structure Analysis**:} \\
\texttt{    *   \$gyzp \textbackslash{}rightarrow jyka\$} \\
\texttt{    *   \$gyzp \textbackslash{}rightarrow lirg\$} \\
\texttt{    *   \$jyka \textbackslash{}rightarrow lirg\$} \\
\texttt{    *   In this structure, \$gyzp\$ acts as a **common cause** (confounder) for both \$jyka\$ and \$lirg\$. This creates a backdoor path from \$jyka\$ to \$lirg\$ via \$gyzp\$.} \\
\\ 
\texttt{2.  **Evaluation of Method 2 (General Correlation)**:} \\
\texttt{    *   Method 2 looks at the raw correlation between \$jyka\$ and \$lirg\$ without conditioning on \$gyzp\$.} \\
\texttt{    *   Because \$gyzp\$ causes both \$jyka\$ and \$lirg\$, there is a spurious association (confounding bias) between them. Even if \$jyka\$ had no causal effect on \$lirg\$, the correlation would exist solely due to the shared cause \$gyzp\$.} \\
\texttt{    *   Therefore, Method 2 cannot isolate the direct causal effect of \$jyka\$ on \$lirg\$ because it fails to control for the confounder \$gyzp\$.} \\
\\ 
\texttt{3.  **Evaluation of Method 1 (Conditional Correlation)**:} \\
\texttt{    *   Method 1 looks at the correlation between \$jyka\$ and \$lirg\$ *case by case according to* \$gyzp\$. This is equivalent to conditioning on \$gyzp\$ (stratifying by levels of \$gyzp\$).} \\
\texttt{    *   By fixing the value of the confounder \$gyzp\$, we block the backdoor path from \$jyka\$ to \$lirg\$ that runs through \$gyzp\$.} \\
\texttt{    *   Once \$gyzp\$ is controlled for, the remaining correlation between \$jyka\$ and \$lirg\$ reflects only the direct causal effect of \$jyka\$ on \$lirg\$.} \\
\\ 
\texttt{**Conclusion**:} \\
\texttt{Since \$gyzp\$ is a confounder that creates a spurious association between \$jyka\$ and \$lirg\$, the general correlation (Method 2) is biased. Conditioning on the confounder (Method 1) is the standard causal inference procedure (specifically, closing the backdoor path) to obtain an unbiased estimate of the causal effect. Thus, Method 1 is more correct.} \\
            }
        }
    }
    & & \\ \\

    \theutterance \stepcounter{utterance}  
    & & & \multicolumn{2}{p{0.3\linewidth}}{
        \cellcolor[rgb]{0.9,0.9,0.9}{
            \makecell[{{p{\linewidth}}}]{
                \texttt{\tiny{[GM$|$GM]}}
                \texttt{Yes} \\
            }
        }
    }
    & & \\ \\

    \theutterance \stepcounter{utterance}  
    & & & \multicolumn{2}{p{0.3\linewidth}}{
        \cellcolor[rgb]{0.9,0.9,0.9}{
            \makecell[{{p{\linewidth}}}]{
                \texttt{\tiny{[GM$|$GM]}}
                \texttt{yes} \\
            }
        }
    }
    & & \\ \\

    \theutterance \stepcounter{utterance}  
    & & & \multicolumn{2}{p{0.3\linewidth}}{
        \cellcolor[rgb]{0.9,0.9,0.9}{
            \makecell[{{p{\linewidth}}}]{
                \texttt{\tiny{[GM$|$GM]}}
                \texttt{game\_result = WIN} \\
            }
        }
    }
    & & \\ \\

\end{supertabular}
}

\end{document}
